\section{Introduction}

This chapter presents the empirical results of PatchOps. The results are organized
around four headline findings rather than around a chronological list of phases. This structure matters because
the project includes both positive gains and honest negative results. The goal of the chapter is therefore not
to maximize apparent success, but to show what the experiments establish about diagnosis, localization, repair,
safety, abstention, and patch honesty.

Every quantitative claim in this chapter is grounded in the canonical per-phase \texttt{findings.md} files. When
summary documents and phase findings differ, the phase findings govern the wording used here. This is
particularly important for the localization studies and for the agentic-abstention result, where the correct
claim depends on reconciled runs rather than on an early single-run summary.

\section{Diagnosis Quality and Early Benchmark Behavior}

Before PatchOps can repair or abstain intelligently, it has to classify failures reasonably well. The scaled
diagnosis ablation ladder provides the clearest diagnosis snapshot on a stratified CI-Repair-Bench subset
(\(n=30\)) \cite{cirepairbench2026}. In that study, category accuracy was 0.83 for the cleaned-log condition and rose to 0.87 when the
workflow YAML was added. Importantly, this gain came from workflow context rather than from deeper repository
context. The full repository-aware configuration did not improve category accuracy beyond the workflow-aware
setting on that slice.

This result is modest rather than dramatic, but it matters. It shows that the diagnosis layer is reasonably
competent and that workflow context is useful, while also suggesting that the primary bottleneck of PatchOps is
not raw failure categorization. This interpretation is consistent with the overall project trajectory: category
accuracy becomes usable relatively early, while later phases reveal that localization and semantic repair remain
substantially harder.

The same diagnosis ablation also exposes category asymmetry. PatchOps is strong on code linting, formatting,
dependency issues, runtime errors, package installation, environment errors, assertion errors,
documentation/docstring failures, and type checking on that slice. It is weaker on test failures and especially
weak on configuration errors. This category split is important because it anticipates the later repair results:
mechanical or localized failures are often understandable, while configuration-heavy and semantically ambiguous
failures are much harder to route cleanly.

\begin{figure}[!htbp]
\centering
\includegraphics[width=0.86\textwidth]{F5_ablation_ladder.png}
\caption{Diagnosis ablation ladder showing the effect of context and gating on diagnosis and localization.}
\label{fig:ablation-ladder}
\end{figure}

\begin{figure}[!htbp]
\centering
\includegraphics[width=0.86\textwidth]{F7_per_category.png}
\caption{Per-category diagnosis behavior on the scaled benchmark slice.}
\label{fig:per-category}
\end{figure}

\section{Fault Localization Remains Hard Across Tested Methods}

The strongest multi-study result in the thesis is that fault localization remains the central bottleneck and
that no tested method clears the target bar. This finding emerges cumulatively across the log-only,
whole-repository, semantic-embedding, and bounded-agentic localization studies.

\subsection{Log-Only Localization Strategies}

The first direct localization study compared four strategies on a stratified CI-Repair-Bench subset (\(n=50\)):
log-path ranking, category reranking, one-shot LLM prediction, and a hybrid strategy. The best result came from
the corrected LLM-first hybrid, which reached Recall@3 of 0.32, Recall@1 of 0.27, and Hit@5 of 0.50. The LLM
was the strongest single signal, while the deterministic methods remained cheaper but less complete.

This was a real improvement over the earlier localization red flag around 0.29, but it was also an honest
warning. Even the best strategy was still far below the advisor's target bar of 0.45. The phase therefore
established two things simultaneously: localization can be improved by better ranking strategy, but the task is
still genuinely difficult.

\begin{figure}[!htbp]
\centering
\includegraphics[width=0.86\textwidth]{F8_localization_recall.png}
\caption{Log-only localization strategies. The corrected LLM-first hybrid is strongest but remains below the target bar.}
\label{fig:phase6-localization}
\end{figure}

\subsection{Whole-Repository Retrieval}

We next tested whether whole-repository retrieval could break that ceiling. In the
deterministic full-repository study (\(n=30\), coverage 1.00), the fused \texttt{full\_repo} strategy became the
best zero-cost retrieval method, reaching Recall@3 of 0.25, Hit@3 of 0.37, and Hit@5 of 0.40. This improved on
the best deterministic log-only baseline, but it still did not approach the 0.45 target. A particularly useful
sub-finding was that pure BM25 over the whole repository was actually harmful on its own, collapsing to
Recall@3 of 0.06. The retrieval signal only became useful when fused with the precise log-path information.

The same phase also measured a paid LLM lever on top of the retrieval stack. The LLM hybrid reproduced the
Recall@3 0.32 result from the earlier log-only study, but an LLM rerank over the BM25 pool did not help whole-repo retrieval. In
fact, it reduced the \texttt{full\_repo} result from Recall@3 0.25 to 0.21. This is one of the most instructive
negative results in the project: adding an LLM on top of lexical retrieval did not solve the localization
ceiling.

\subsection{Semantic Embedding Retrieval}

We then tested the remaining obvious retrieval lever: semantic embeddings. The result was another honest
negative. On the same \(n=30\) draw, \texttt{full\_repo\_embed} reached Recall@3 0.21, which did not beat the
deterministic BM25 fusion at 0.25. It traded a little top-3 precision for slightly better depth at rank 5, but
it still failed to cross the target bar.

By the end of this retrieval comparison, PatchOps had measured three distinct retrieval families on this task:

\begin{enumerate}
\item lexical whole-repository retrieval, peaking at Recall@3 0.25,
\item one-shot LLM reranking or hybridization, peaking around Recall@3 0.27--0.32, and
\item semantic embedding reranking, at Recall@3 0.21.
\end{enumerate}

All three remained below 0.45. This is the point at which the thesis can defend a scoped interpretation:
localization in this setting is not merely solved by swapping retrieval methods. Across the methods tested here,
simple retrieval is insufficient, and the residual difficulty appears to involve repository reasoning.

\subsection{Bounded Agentic Exploration}

The final major localization lever was bounded agentic exploration over a repository clone. This is
where the reconciliation note matters most. The canonical result for the gpt-4o-mini agent on the \(n=30\)
seed-7 slice is Recall@3 0.24, with average latency 19.9 seconds per instance and a high forced-fallback rate.
The agent therefore slightly edges the deterministic full-repository baseline but does not beat the LLM hybrid
and remains far below 0.45.

This result is powerful precisely because it is the fourth negative. After lexical retrieval, LLM one-shot
reranking, semantic embeddings, and now agentic exploration, localization still does not cross the bar. The
task appears bottlenecked by repository reasoning difficulty in this evaluated setting rather than by the
absence of one more retrieval heuristic.

There is one important nuance. The agent does occasionally show the \emph{possibility} of developer-style
exploration helping where ranking fails. On the few instances where it actually commits to a \texttt{finish}
action rather than timing out into fallback, its committed subset looks stronger than the deterministic baseline,
and at least one clean win exists where repository exploration found a relevant fix file the retrieval baseline
missed. However, these wins are too sparse and unstable to overturn the global conclusion.

\subsection{Localization Conclusion}

Taken together, these four localization studies support the first headline finding:
\emph{across the four methods tested, simple retrieval is insufficient in this setting, and the residual
difficulty appears to be repository reasoning.} The best measured values across these families remain clustered
well below the target threshold. This does not mean retrieval is useless, and it is not a universal law about
all possible localizers. It means that, for PatchOps and this benchmark setting, retrieval is necessary but not
sufficient, and the downstream system must remain honest when localization is weak.

\begin{figure}[!htbp]
\centering
\includegraphics[width=0.88\textwidth]{F11_localization_strategies.png}
\caption{Localization strategies across Phases 6, 8, and 10. None clears the 0.45 Recall@3 target.}
\label{fig:localization-strategies}
\end{figure}

\begin{figure}[!htbp]
\centering
\includegraphics[width=0.88\textwidth]{F18_agent_localization.png}
\caption{Agentic localization result and termination behavior. Agentic exploration is useful in isolated cases but does not clear the localization plateau.}
\label{fig:agent-localization}
\end{figure}

\begin{table}[!htbp]
\centering
\small
\caption{Headline localization results for the localization bottleneck study.}
\label{tab:localization-headline}
\begin{tabularx}{\textwidth}{
  >{\raggedright\arraybackslash}p{0.29\textwidth}
  c
  c
  >{\raggedright\arraybackslash}X
}
\toprule
Method family & Sample & Recall@3 & Interpretation \\
\midrule
LLM-first hybrid & \(n=50\) & 0.32 & Best overall measured localizer. \\
Full-repository fused BM25 & \(n=30\) & 0.25 & Best zero-cost full-repo retrieval. \\
Embedding rerank & \(n=30\) & 0.21 & Does not beat lexical retrieval. \\
Bounded agentic exploration & \(n=30\) & 0.23--0.24 & Does not beat the best hybrid. \\
\bottomrule
\end{tabularx}
\end{table}

\noindent\textit{Comparability note.} The best hybrid result is measured on a 50-instance slice, whereas the
whole-repository, embedding, and agentic rows are measured on a 30-instance slice. The fine-grained ordering
between these families should therefore be read as indicative, while the robust conclusion is the shared
plateau below the 0.45 Recall@3 target.

\section{Better Localization Is Necessary but Not Sufficient for Repair}

The second headline finding comes from the localization-to-repair linkage experiment. Here the
methodology asked a decisive question: if localization improves and patches become cleaner, does CI Pass@1 rise
with it?

On the reproducible subset (\(n=24\)) of the scaled repair re-run, the ungated \texttt{raw\_llm} baseline reached
Pass@1 0.12 with apply-rate 0.43. The gated \texttt{repo} variant, which used oracle changed-files context from
the earlier setup, had coverage 0.29, apply-rate 0.29 among attempted cases, and Pass@1 0.00. The key result
was the \texttt{repo\_full} variant: once real clone-based full-repository retrieval was used, apply-rate among
attempted cases rose to 1.00, yet Pass@1 stayed 0.00.

These two Pass@1 values are not directly comparable across variants: \texttt{raw\_llm} is the ungated baseline
at 0.96 coverage, whereas \texttt{repo\_full} is the gated path at 0.29 coverage. The finding rests on the
apply-rate-to-Pass@1 decoupling within the gated path, not on a direct cross-variant Pass@1 comparison.

This is a central thesis result because it isolates a specific bottleneck transition. Earlier, one could have
suspected that repair was failing mainly because the generated diffs did not even apply cleanly. The linkage study shows
that once the apply-rate bottleneck is removed, the system still fails to make CI green on the gated path. The
remaining problem is semantic correctness, not simply file selection or diff validity.

In other words, localization improvement is necessary because it prevents invalid or off-tree patches, but it is
not sufficient to produce correct repair. This is exactly why the thesis emphasizes honest abstention and safety
controls. When the system reaches the semantic frontier, caution becomes more valuable than overconfident
automation.

\begin{figure}[!htbp]
\centering
\includegraphics[width=0.70\textwidth]{F12_apply_bottleneck.png}
\caption{Localization-to-repair linkage: full-repository retrieval removes the apply-rate bottleneck among attempted patches, but repair success remains blocked by semantic correctness.}
\label{fig:apply-bottleneck}
\end{figure}

\begin{figure}[!htbp]
\centering
\includegraphics[width=0.82\textwidth]{F9_repair_passk.png}
\caption{Docker repair evaluation on the reproducible subset. Ungated repair obtains some CI-green outcomes, while the gated policy remains conservative.}
\label{fig:repair-passk}
\end{figure}

\begin{table}[!htbp]
\centering
\small
\caption{Localization-to-repair linkage on the reproducible Docker subset.}
\label{tab:repair-linkage}
\begin{tabular}{lccc}
\toprule
Variant & Coverage & Apply-rate among attempted & CI Pass@1 \\
\midrule
\texttt{raw\_llm} & 0.96 & 0.43 & 0.12 \\
\texttt{repo} & 0.29 & 0.29 & 0.00 \\
\texttt{repo\_full} & 0.29 & 1.00 & 0.00 \\
\bottomrule
\end{tabular}
\end{table}

\subsection{Supporting Result: Learned Classification as a Reproducibility and Cost Result}

The learned failure classifier adds a narrower but still important result. On a held-out test split of
\(n=118\), the learned model reached exact accuracy 0.74, lenient accuracy 0.84, and macro-F1 0.58. The rule
baseline, by contrast, achieved exact accuracy 0.37 and lenient accuracy 0.47. The learned model therefore
nearly doubled the deterministic rule baseline and landed inside the LLM classifier's observed lenient accuracy
band.

This result should be interpreted carefully. The classifier is not a repair-success breakthrough. Its
value is that it makes one part of the pipeline cheaper, more reproducible, and locally trainable while keeping
accuracy competitive with the LLM path on the categories that dominate the dataset. This is why it belongs in
the results chapter but not as the main headline. It strengthens the system scientifically without pretending to
solve the core repair problem.

\begin{figure}[!htbp]
\centering
\includegraphics[width=0.82\textwidth]{F13_classifier.png}
\caption{Learned classifier compared with the rule baseline. The classifier is a reproducibility and cost contribution, not a repair-rate claim.}
\label{fig:classifier}
\end{figure}

\section{CI-Green Is Not Enough}

The third headline finding is that a green or plausible-looking result is not a sufficient indicator of
correctness. PatchOps investigates this in two complementary ways: reward-hacking detection and patch
plausibility analysis.

\subsection{Reward Hacking}

The reward-hacking audit asked whether any of the 66 produced patches achieved apparent progress by cheating rather
than by fixing the underlying issue. The result was small but important: the observed cheat-rate was non-zero,
with one detected reward-hacking patch among the 66 produced patches. That patch deleted a test and removed an
assertion. The detector caught it, and the wired gate would block it.

Equally important, none of the three CI-green patches in the run were hacky. This lets the thesis state the
result honestly: reward hacking is not the dominant behavior of this model on this workload, but it is real
enough to require explicit defenses. This is exactly the type of low-frequency, high-consequence behavior a
safety-first CI agent should guard against.

\begin{figure}[!htbp]
\centering
\includegraphics[width=0.82\textwidth]{F16_reward_hacking.png}
\caption{Reward-hacking audit. The observed cheat rate is low but non-zero, justifying an explicit detector and gate.}
\label{fig:reward-hacking}
\end{figure}

\subsection{Patch Plausibility}

The plausibility study asked a different but related question: even when a patch looks syntactically reasonable and passes the
safety gate, is it editing a file that the developer's actual fix touched? On the 66 produced patches in the
validation set, the base on-target rate was only 0.14. In other words, most generated patches were not even
editing a gold-overlapping file.

The static plausibility score performed well as a discriminator, reaching AUC 0.77 overall, with evidence
locality alone already at 0.75. When used as a gate, it materially improved auto-PR precision: from 0.14 at
full coverage to 0.27 at 50\% coverage, 0.35 at 25\% coverage, and a peak around 0.50 at very low coverage.

This is exactly the patch-level counterpart of the selective-prediction story from the abstention studies. At the diagnosis
level, calibrated confidence tells PatchOps when to act. At the patch level, plausibility tells PatchOps when a
patch should be previewed instead of being treated as an automatic action candidate.

\begin{figure}[!htbp]
\centering
\includegraphics[width=0.82\textwidth]{F17_patch_plausibility.png}
\caption{Plausibility gate. Precision improves as automatic action is restricted to more plausible patches.}
\label{fig:patch-plausibility}
\end{figure}

\begin{table}[!htbp]
\centering
\small
\caption{Patch honesty and plausibility results.}
\label{tab:honesty-plausibility}
\begin{tabularx}{\textwidth}{
  >{\raggedright\arraybackslash}p{0.28\textwidth}
  c
  >{\raggedright\arraybackslash}X
}
\toprule
Result & Value & Interpretation \\
\midrule
Reward-hack patches & 1/66 & Low-frequency but real cheating behavior. \\
CI-green patches that were hacky & 0/3 & No observed green patch was gamed in this run. \\
Base on-target rate & 0.14 & Most produced patches do not edit a gold-overlapping file. \\
Combined plausibility AUC & 0.77 & Static signals predict on-target behavior. \\
Peak gated precision & \(\approx 0.50\) & Higher precision is possible at low coverage. \\
\bottomrule
\end{tabularx}
\end{table}

\section{Abstention Must Be Engineered}

The fourth headline finding emerges most clearly in the selective prediction and safety results. The thesis
claim is not simply that PatchOps abstains. The claim is that \emph{good abstention must be engineered through
external controls rather than left to the free-form model alone.}

\subsection{Selective Prediction and Calibration}

The selective-prediction study turns abstention into a rigorous operating characteristic. On the held-out classification split
(\(n=118\)), the learned classifier's confidence produced a strong risk-coverage curve: base accuracy 0.84,
AURC 0.051, and selective accuracy rising to 0.97 at 50\% coverage and 1.00 at 25\% coverage. By contrast, the
current heuristic engine confidence was a poor abstention signal, with base accuracy 0.74, AURC 0.261, and
erratic selective behavior.

This is a major thesis-quality result. It shows that PatchOps can abstain intelligently when the right
confidence signal is used, and that the current hand-rolled confidence should not be trusted as the main gating
input. The same study also shows that confidence calibration matters. The learned classifier's raw expected
calibration error of 0.224 was cut to 0.119 by isotonic calibration, making the emitted confidence easier to
use as a stable decision threshold.

\begin{figure}[!htbp]
\centering
\includegraphics[width=0.82\textwidth]{F14_risk_coverage.png}
\caption{Risk-coverage behavior. Learned classifier confidence is a much stronger abstention signal than hand-rolled engine confidence.}
\label{fig:risk-coverage}
\end{figure}

\begin{figure}[!htbp]
\centering
\includegraphics[width=0.62\textwidth]{F15_reliability.png}
\caption{Reliability diagram showing the value of isotonic calibration.}
\label{fig:reliability}
\end{figure}

\subsection{The Learned Calibrated Decision Gate}

The learned decision-gate study turns the earlier calibration lesson into an implemented decision policy. Earlier PatchOps versions
used hand-set constants, including a low confidence floor and a high-confidence patch threshold, to decide when
a low-risk case should act or hold. The learned gate replaces the low-risk act-versus-hold threshold with a learned,
isotonic-calibrated decision gate over pre-generation signals: calibrated classifier confidence, category risk,
flaky status, semantic ambiguity, log-named-file count, and evidence length.

On the expanded Safety Set (\(n=76\)), the learned gate reaches act/hold accuracy 0.91 with a Wilson interval
of [0.82, 0.95] under leave-one-out evaluation and 0.92 under five-fold evaluation, compared with 0.855 for the
hand-tuned threshold policy. It is also well calibrated, with \(P(\mathrm{act})\) ECE approximately 0.07. The
honest interpretation is matches-or-beats rather than a statistically separated win, because the confidence
intervals overlap at this sample size. The important engineering result is that the decision is no longer
explained by hand-set thresholds alone.

The most concrete gain appears on the low-risk lane where the engine actually uses the learned gate. Among the
15 cases routed to that lane, all 15 are true-patch cases. The learned gate acts on 15/15, while the hand-tuned
thresholds act on only 11/15, recovering four confident patches that the older policy held back, with 0 over-act
on that lane. This is a decision-quality result, not a Pass@1 result: it improves the act-versus-hold policy
without claiming that semantic repair success is solved.

This result should also be interpreted with its small trade-offs visible. Under leave-one-out evaluation, the
learned gate incurs a marginal abstention and over-act cost relative to the hand-tuned policy
\((0.95 \text{ vs. } 0.96;\ 3/56 \text{ vs. } 2/56)\), but it gains higher overall act/hold accuracy and
recovers low-risk actions that the earlier policy held back.

The same refinement also ships the earlier isotonic calibrator inside the classifier artifact. The emitted classifier
confidence ECE falls from 0.18 to 0.06 in the held-out check, making calibrated confidence part of the engine
rather than only an offline analysis. A related embedding-versus-TF-IDF classifier ablation is an honest null:
neural embeddings reach lenient accuracy 0.66, below the shipped TF-IDF classifier at 0.83. This reinforces the
earlier interpretation that classification is not the bottleneck and that simple lexical signals remain strong
for this stage.

The two calibration figures in this section refer to different evaluations: the earlier selective-prediction
benchmark reports isotonic calibration reducing ECE from 0.224 to 0.119, while the learned decision-gate study reports the
shipped classifier artifact on the later held-out check, where ECE falls from 0.18 to 0.06. They are therefore
consistent calibration results, not competing measurements of the same run.

\begin{figure}[!htbp]
\centering
\includegraphics[width=0.86\textwidth]{F19_decision_gate.png}
\caption{Learned decision gate. The calibrated policy matches-or-beats the hand-tuned threshold policy and recovers held-back low-risk actions without over-acting on that lane.}
\label{fig:decision-gate}
\end{figure}

\begin{table}[!htbp]
\centering
\small
\caption{Learned decision gate summary.}
\label{tab:decision-gate}
\resizebox{\textwidth}{!}{%
\begin{tabular}{lccccc}
\toprule
Policy & Act/hold acc. & Abstention acc. & Over-act & \(P(\mathrm{act})\) ECE & AURC \\
\midrule
Hand-tuned thresholds & 0.855 & 0.96 & 2/56 & -- & -- \\
Learned gate, LOO & 0.908 [0.82, 0.95] & 0.95 & 3/56 & 0.072 & 0.106 \\
Learned gate, 5-fold & 0.921 [0.84, 0.96] & 0.96 & 2/56 & 0.066 & 0.089 \\
\bottomrule
\end{tabular}
}
\end{table}

\subsection{Safety Set Evaluation}

The earlier safety-set study already demonstrated the operational value of abstention-oriented behavior. On
the labeled \(n=20\) safety benchmark, PatchOps achieved abstention accuracy 0.92, flaky F1 0.86, and 0/20
unsafe auto-fixes. It also correctly abstained on all escalate and rerun cases, with one over-eager preview
versus patch mismatch that was later addressed by the semantic ambiguity backstop.

Methodologically, the most important part of this result is not that the benchmark is small, but that the
decision classes are explicit. The system is being evaluated on whether it chooses the right kind of action, not
only on whether it emits diffs. The Safety Set is later expanded to \(n=76\) for the learned decision-gate study, but the original
numbers remain the original 20-case result and should be read with that denominator.

\begin{figure}[!htbp]
\centering
\includegraphics[width=0.78\textwidth]{F6_safety_gate.png}
\caption{Safety Set behavior. The original 20-case result measures decision accuracy, abstention, and unsafe auto-fix behavior.}
\label{fig:safety-gate}
\end{figure}

\subsection{Self-Abstention Nuance}

The agentic localization phase supplies the most important cautionary nuance in the whole thesis. The gpt-4o-mini
agent was explicitly given an \texttt{abstain} action, yet across two runs on the same slice it abstained only
rarely and run-dependently: 0/30 in one run and 2/30 in the canonical reconciled run. When those abstentions
did occur, they were justified. However, the correct conclusion is not that the model abstains well. The
correct conclusion is that self-abstention is too rare and unstable to rely on.

This is why the wording matters. The robust thesis claim is that the agent abstains \emph{rarely and
run-dependently}, not never. That instability is itself a finding. It shows that if PatchOps is to abstain
reliably, the abstention logic has to be engineered into the decision stack through calibrated confidence and
external safety controls. The learned decision-gate result is the capstone of this argument: the act-versus-hold
decision is explicitly learned and calibrated rather than left to self-abstention or hand-set thresholds alone.

\section{Safety and Practicality as System Outcomes}

The combined results also say something broader about the system's practical operating profile. On the safety
benchmarks, PatchOps is conservative and generally asymmetric toward caution. On the repair benchmarks, the
ungated baseline can produce more aggressive behavior and occasionally more green outcomes, but at the cost of
weak guarantees and lower interpretability. The gated system is more trustworthy, but this trustworthiness is
paid for with lower coverage.

This trade-off is not a flaw in the thesis. It is part of the thesis claim. PatchOps is not trying to maximize
raw patch volume. It is trying to define a repair agent that behaves sensibly when confidence and correctness
signals are weak.

The cost numbers also support this practical framing. The scaled diagnosis ablation remained in the low-dollar
range, with approximately \$1.47 total spend for 120 diagnoses using the cheaper model configuration. The later
whole-repository retrieval and learned-classifier studies together cost approximately \$0.57. The agentic localization study remained cheap in
absolute dollar terms but expensive in latency, averaging about 20 seconds per instance and delivering no Recall@3
gain over the best one-shot hybrid localizer. This makes it clear that agentic exploration is not currently the
default practical choice.

\begin{figure}[!htbp]
\centering
\includegraphics[width=0.78\textwidth]{F3_token_cost.png}
\caption{Token-cost profile for the scaled evaluation configuration.}
\label{fig:token-cost}
\end{figure}

\section{Result Synthesis}

The results chapter supports the thesis spine in a coherent way.

First, diagnosis quality becomes serviceable relatively early, especially when workflow context is added, but it
does not remove the deeper localization and repair bottlenecks.

Second, localization remains the defining challenge, and repeated attempts to solve it through new retrieval
mechanisms and bounded agentic exploration plateau below the desired threshold. This is the strongest evidence
that, in the evaluated PatchOps setting, the unresolved part of CI repair appears to involve repository
reasoning rather than simple matching alone.

Third, once localization is improved enough to make patches apply cleanly, the next wall is semantic
correctness. That transition is one of the clearest and most thesis-worthy findings in the project.

Fourth, confidence and safety controls matter because the system cannot rely on the base agent to abstain
reliably, and because CI-green or plausible-looking patches are not trustworthy enough on their own.
The learned decision gate strengthens this point by replacing the old low-risk confidence threshold with a learned calibrated
gate that matches-or-beats the threshold baseline while recovering held-back low-risk actions.

\section{Chapter Conclusion}

This chapter presented the empirical backbone of the PatchOps thesis. The results show that workflow context
improves diagnosis modestly, that localization remains the central bottleneck across four different method
families, that fixing localization is necessary but not sufficient for repair success, that a learned classifier
can deliver strong reproducibility and cost benefits, that abstention becomes useful only when engineered
through calibrated external controls, that the learned decision-gate study turns the act-versus-hold decision into a learned calibrated
gate, and that CI-green is too weak a success criterion without reward-hack and plausibility defenses.

Taken together, these results justify the central thesis claim: PatchOps is best understood not as a system that
``solves'' CI repair, but as a reproducible, safety-gated, abstention-aware repair agent that remains honest
about where repository reasoning and semantic repair are still hard. The next chapter synthesizes these results
into the broader discussion and implications of the work.
